\documentclass[a4paper]{article}
\usepackage[pdftex]{graphicx}
\usepackage[utf8]{inputenc}
\usepackage{enumerate}
\usepackage{amssymb}
\usepackage{color}
\usepackage{tikz}
\usepackage{href-ul}
\hypersetup{
	colorlinks=true,
	linkcolor=blue,
	urlcolor=blue}
\usepackage{geometry}
\geometry{a4paper, top=15mm, left=15mm, right=15mm, bottom=15mm,
	headsep=10mm, footskip=12mm}

\begin{document}
\begin{tikzpicture}(10,3)
	\draw[ultra thick](2,0) --(10,0) -- (10,1.3) --(2,1.3) -- (2,0);
	\draw[fill=black](2,0)-- (0,.6) -- (2,.6) -- (2,0);
	\node at (6,.6) {\large Lösungsblatt von \href{https://www.okuyakl.de}{www.okuyakl.de}};
\end{tikzpicture}
\vspace{0.5 cm}

\noindent{\bf Aufgabe 1.}\\
$${1 \over 2}\cdot (8x + y)^2 = 32x^2 + 8xy + {1\over 2} y^2$$

\noindent{\bf Aufgabe 2. a)}\\
Wir formen um: 
$$-4x^2 + 2x +2 \ge 0 $$
Zum Faktorisieren setzen wir in die Mitternachtsformel ein:
$$ x_{1/2} = { -2 \pm \sqrt{2^2 - 4\cdot (-4)\cdot 2} \over 2\cdot (-4)} = {-2 \pm 6 \over -8} $$
$$ x_1 = -{1\over 2} \qquad x_2= 1$$
Mit den erhaltenen Lösungen und dem Vorfaktor $a=-4$ formen wir um zu:
$$-4\left(x+{1\over 2}\right)(x-1)\ge 0 \qquad (*)$$
Für dieses Produkt erstellen wir eine Vorzeichentabelle:

$$
\renewcommand{\arraystretch}{2}
\begin{array}{r|c|c|c|c|c}
& x < -{1\over 2} & x =-{1\over 2} & -{1\over 2} < x < 1 & x=1 & x>1 \\
\hline
-4 & - & - & - & - & - \\
\left(x+{1\over 2}\right) & - & 0 & + & + & + \\ 
(x-1) & - & - & - & 0 & + \\
\hline
(*) & - & 0 & + & 0 & - 
\end{array}
$$
Es folgt für die Lösungsmenge:
$\mathbb{L}=\{x|-{1\over 2} \le x \le 1\} $\\
\setlength{\unitlength}{1 cm}
\begin{picture}(10,2)
\put(0,1){\vector(1,0){9}}
\put(0,0.85){\line(0,1){0.3}}
\put(2,0.85){\line(0,1){0.3}}
\put(4,0.85){\line(0,1){0.3}}
\put(6,0.85){\line(0,1){0.3}}
\put(8,0.85){\line(0,1){0.3}}
\put(-0.2,0){-1}
\put(1.6,0){$-{1\over 2}$}
\put(3.9,0){0}
\put(5.9,0){${1\over 2}$}
\put(7.9,0){1}
\put(1.93,0.9){$\Big[$}
\put(7.93,0.9){$\Big]$}
\linethickness{2 pt}
\color{green}
\put(1.9,1){\line(1,0){6}}
\end{picture}

\noindent {\bf Aufgabe 2. b)}\\
Wir gehen analog zu a) vor:
$$ - {1 \over 2} x^2 + 2x -2 > 0 $$
Die Mitternachtsformel liefert:
$$x_{1/2}= { -2 \pm \sqrt{4- 4\cdot \left(-{1 \over 2}\right)\cdot (-2)} \over 2 \cdot \left(-{1 \over 2}\right)}= {-2 \pm 0 \over -2} = 1 $$
Es gibt nur eine Lösung. Wir faktorisieren...
$$-{1 \over 2}\cdot(x-1)^2>0 \qquad (**)$$  
...und erstellen eine Vorzeichentabelle:
$$
\renewcommand{\arraystretch}{2}
\begin{array}{r|c|c|c}
& x < 1 & x = 1 &  x>1 \\
\hline
-{1 \over 2} & - & - & - \\
(x-1)^2 & + & 0 & +  \\ 
\hline
(*) & - & 0 & - 
\end{array}
$$
Die Lösungsmenge ist, weil es keine Werte größer als null gibt, leer:
$\mathbb{L}=\{ \}$

\noindent {\bf Aufgabe 3.}\\
Dies ist nur eine quadratische Gleichung, wenn $a \ne 0 $

Damit erhalten wir für Fall 1: $a=0 \Rightarrow$ lineare Gleichung $\Rightarrow$ Eine Lösung 

Fall 2: $a \ne 0 $
Zur Bestimmung der Anzahl der Lösungen bestimmen wir die Diskriminante $b^2-4ac$:
$$D=\left(-{3 \over 2}\right)^2-4\cdot a \cdot 1 = {9\over 4} - 4a$$
Wir setzen nun $D=0$ und erhalten für $a$:
$$
\renewcommand{\arraystretch}{2}
\begin{array}{rcll}
{9\over 4} - 4a & = & 0 & | + 4a \\ 
	{9\over 4}  & = & 4a & | :4 \\
	{9\over 16} & = & a 
\end{array}
$$
Zur weiteren Fallunterscheidung erstellen wir eine Vorzeichentabelle:
$$
\renewcommand{\arraystretch}{2}
\begin{array}{r|c|c|c}
& a < {9\over 16} & a = {9\over 16} &  a>{9\over 16} \\
\hline
{9\over 4} - 4a & + & 0 & - \\
\hline
D & + & 0 & -\\
\hline
\textnormal{Anzahl Lösungen} & 2 & 1 & 0 
\end{array}
$$

Für $a \in ] - \infty ; 0 [ \quad \vee \quad ]0;{9\over 16}[ $ gibt es zwei Lösungen;
Für $a={9\over 16}$ oder $a=0$ gibt es eine Lösung und für $a>{9\over 16}$ gibt es keine Lösung. 

\noindent {\bf Aufgabe 4.}\\
Wir setzen Parabel- und Geradenterm gleich und vereinfachen:
$$
\begin{array}{rcll}
mx + {3 \over 2} & = & x^2 - {3 \over 2}x + {3 \over 2} & |-{3 \over 2}\\
          mx & = & x^2 - {3 \over 2}x & | - mx \\
          0  & = & x^2 - \left({3 \over 2} + m\right)x \\
          0  & = & x \cdot \left(x - {3 \over 2}-m \right)    
\end{array}
$$
Ein Produkt ist gleich 0, wenn einer der Faktoren null wird; daraus folgt:
$$x_1=0$$
$$x_2 - {3 \over 2}-m =0$$
$$x_2 = {3 \over 2} + m $$

Wir haben nun zwei Lösungen für den Schnitt der Geraden mit der Parabel. Die Gerade soll jedoch {\it Tangente} sein, und diese Tangente hat nur {\it einen}  Berührpunkt. Dies bedeutet, dass wir $m$ so wählen, dass beide Lösungen für $x_1$ und $x_2$ identisch sind,
also 
$$x_2=0 \quad \Rightarrow \quad m= -{3 \over 2}$$
Skizze für $m=1$ 
\includegraphics[width=5 cm]{paga157}

\fbox{
	\begin{minipage}{0.48\textwidth}
		\noindent{\bf Aufgabe 5. a)}\\
		Zum Berechnen der Schnittpunkte setzen wir beide Terme gleich. Zuerst müssen wir die y-Seiten gleich machen. Hier bietet sich an, den ersten Term mit $(-4)$ zu multiplizieren. 
		$$
		\renewcommand{\arraystretch}{2}
		\begin{array}{rrcll}
		I:& -4y &=& -{9 \over 4}x + 5 \\
		II:& -4y &=& -ax -5 \\
		I=II & -{9 \over 4}x + 5 &=& -ax -5 &| -5 + ax \\
		& -{9 \over 4}x + ax &=&  -10 \\
		&  (-{9 \over 4} + a)\cdot x &=&  -10  &|\cdot (-4)\\
		& (9-4a)\cdot x &=& 40 &|:(9-4a)\\
		& x &=& {40 \over 9-4a} \\
		\end{array}
		$$
	\end{minipage}
}
\hfill
\fbox{
	\begin{minipage}{0.48\textwidth}
		\noindent{\bf Aufgabe 5. b)}\\
		Wenn der Nenner des Terms null wird, gibt es keine Lösung und damit keinen Schnittpunkt:
		$$9-4a=0 \quad \Rightarrow \quad a={9\over 4}$$
		In diesem Fall sind die beiden Geraden parallel.	
		
		\noindent{\bf Aufgabe 5. c)}\\
		$$
		\renewcommand{\arraystretch}{2}
		\begin{array}{rcll}
		-8 &=& {40 \over 9-4a} &| \cdot (9-4a)\\
		-72 + 32a &=& 40 &|+72\\
		32a &=& 112 &|:32 \\
		a &=& 3,5
		\end{array}
		$$	
	\end{minipage}
}

\begin{center}
	\includegraphics[width=7 cm]{../../viecher/endcomic.pdf}
	
	Hier geht es zurück zum \href{https://www.okuyakl.de/math/m9quglL040/aa040.pdf}{Aufgabenblatt}
\end{center}

\end{document}

